\section{Experimental Design}
\label{sec:design}

The evaluation separates the questions of threshold adaptation, batch interpretation, domain-rule coverage and execution cost. Each component is compared under a common replay protocol, while the external benchmarks and chronological replay test aspects that controlled injection cannot establish.

\subsection{Datasets}
\label{sec:design:data}

Three public datasets provide contrasting data distributions and validation requirements (Table~\ref{tab:datasets}). BTS Airline contains US Bureau of Transportation Statistics on-time performance records for January to June 2024~\cite{BTS2024}. Online Retail II contains invoice lines from a UK online gift-ware retailer between December 2009 and December 2011~\cite{Chen2019}. Chicago Food Inspections contains municipal inspection records from 2010 onward~\cite{Chicago2026}. Flight delays and taxi times have long numerical tails; retail quantities and prices are discrete and highly skewed; inspection outcomes are predominantly categorical and include geographic coordinates. These differences allow the same adaptation policy to be tested where its assumptions are more and less appropriate.

For the controlled experiments, preprocessing removed specified natural invalidities before the replay and injection stages. Cancelled and diverted flights, cancelled invoices, non-positive quantities or prices, records with missing monitored values and exact duplicates were removed. This creates a consistent basis for an audited injected-error answer key. It does not establish that every remaining source value is factually correct. Accordingly, controlled precision and recall refer to the logged injected defects under the stated cleaning policy.

\begin{table}[t]
\caption{Datasets used in the controlled experiments. Batch sizes are those of the test split. Numerical and categorical monitored columns are separated by a semicolon.}
\label{tab:datasets}
\small
\begin{tabularx}{\linewidth}{@{}lXrrrX@{}}
\toprule
Dataset & Domain & Raw rows & Clean rows & Batch & Monitored columns \\
\midrule
BTS Airline & Transport & 3,461,319 & 3,403,465 & 50,000 & DepDelay, TaxiOut, AirTime, Distance; Reporting\_Airline, Origin, OriginState \\
Online Retail II & E-commerce & 1,067,371 & 1,007,911 & 30,000 & Quantity, Price; Country, StockCode \\
Chicago Food Inspections & Public health & 315,856 & 309,460 & 10,000 & Latitude, Longitude; Risk, Results, Facility Type \\
\bottomrule
\end{tabularx}
\end{table}

\subsection{Splits, Seeds and Freezing}
\label{sec:design:splits}

A fixed partition seed, independent of the run seeds, divides each cleaned dataset into a 30\% development pool and a 70\% test pool without shared records. Method development, debugging and controlled-experiment LLM prompt examples use the development pool and five development seeds. Values reserved for valid-novelty tests are excluded from prompt examples. The test uses ten paired seeds, 1001 to 1010. The runners refuse to use the test pool until the protocol's frozen flag is enabled.

The internal protocol freeze is dated 20 October 2026, following the fifth complete development run. Run manifests record SHA-256 fingerprints of code, configuration and input data; the principal confirmatory results share code fingerprint \texttt{fd749de031ca4f54}. Development changes are documented with their motivation and observed effects. This provides an auditable record of the implementation and settings used for testing, although it is not independent public preregistration.

Within each run, the chosen record pool is shuffled using the run seed and divided into twenty batches. The first five batches calibrate the monitor, and the remaining fifteen are evaluated. Thus the controlled runs remove natural temporal ordering: the legitimate drift in their answer key is the specified injected change. Section~\ref{sec:design:replay} separately evaluates the unshuffled chronological data. Methods share the same replay and injected records within a seed, allowing paired comparisons.

\subsection{Scenarios}
\label{sec:design:scenarios}

Seventeen scenarios are generated by deterministic in-memory replay (Table~\ref{tab:scenarios}). Their purpose is to vary legitimate change and erroneous values separately before testing their combination.

\emph{Legitimate drift} starts at 2\% in $B_6$ and increases by two percentage points per batch to 30\% in $B_{20}$. A numerical drift column is multiplied by $(1+d)$; a proportion $d$ of a categorical drift column is redrawn from the ten most common categories, each accounting for at least 1\% of the source records. Drift affects delays and taxi times for flights, quantity and price for retail, and risk and result categories for inspections. These choices avoid changes that would be physically inappropriate, such as increasing a route's distance by 30\% or moving all inspection coordinates.

\emph{Background errors} affect 1\% of records in each batch, including warm-up. They are divided equally among missing values, exact duplicate records, invalid categories and numerical outliers placed $U(4,12)$ standard deviations from the batch mean. Duplicates are appended last as copies of untouched records. Labels are audited against the final batch so that an injected operation is not assumed to remain an observable error after later operations. Mild-error scenarios place outliers between three and five standard deviations from the mean.

\emph{Degradation} increases the error rate to 3\% from $B_{11}$, 5\% from $B_{14}$ and 10\% from $B_{17}$. \emph{Contamination} introduces a persistent systematic error in 5\% of records from $B_8$: TaxiOut or Price is multiplied by three. Chicago longitude uses a factor of 1.002, approximately 0.18 degrees, because multiplication by three would produce an obvious geographic violation. Contamination tests whether learning from recurring observations weakens detection.

\emph{Domain-impossible values} affect an additional 1\% of records from $B_6$. These include negative taxi-out times, non-positive airborne times, negative route distances, well-formed but non-existent carrier, airport or state codes, non-positive retail quantities, negative prices, and invalid country, stock, risk, result or facility-type values. They are absent from the cleaned data and warm-up. The selected cases test domain restrictions that a broad statistical range or a syntax pattern may not capture.

\emph{Novel categories} test admission rather than only detection. Real values are withheld from the record pool and introduced at 0.5\% from $B_6$ or $B_{11}$: examples include carriers HA and G4, states HI and AK, airports STL and PDX, countries Australia and Japan, and the inspection result \emph{Not Ready}. Invalid values are introduced at the same rates, including carrier ZX, state QH, misspelled countries and undocumented inspection results. The test contains twelve valid and twelve invalid distinct values. Replaying these values across seeds changes their exposure and recurrence, not the semantic diversity of the test set.

\emph{Incidents} add sparse or temporary defects to the 1\% background: duplicates at 0.5\% from $B_{11}$; invalid categories in a low-cardinality column at 0.5\% from $B_{11}$; duplicates and invalid categories at 3\% in $B_9$ and $B_{10}$; or the same 3\% incident from $B_9$ to $B_{14}$ followed by recovery. These scenarios test whether selection remains responsive when evidence is brief or infrequent.

\begin{table}[t]
\caption{Scenarios. Error rates are shares of records. Gradual drift rises from 2\% at $B_6$ to 30\% at $B_{20}$.}
\label{tab:scenarios}
\small
\begin{tabularx}{\linewidth}{@{}llXl@{}}
\toprule
Scenario & Drift & Errors & Used for \\
\midrule
A stable & none & 1\% background & RQ1, RQ2, RQ3 \\
B drift & gradual & 1\% background & RQ1, RQ2, RQ3 \\
C degradation & none & 1\% to $B_{10}$, then 3, 5, 10\% & RQ1, RQ3 \\
D drift and degradation & gradual & as C & RQ1, RQ3 \\
Z zero-error control & gradual & none & RQ1a \\
E1, E2 mild errors & none; gradual & 1\%, outliers at 3 to 5 sd & RQ1a \\
F1, F2 contamination & none; gradual & 1\% plus 5\% systematic from $B_8$ & RQ1a \\
S1, S2 domain-impossible & none; gradual & 1\% plus 1\% impossible values from $B_6$ & RQ2a, ablation \\
G1, G2 novelty & none; gradual & 1\% plus valid and invalid new categories & RQ2b \\
H1 to H4 incidents & none & 1\% plus sparse, short or recovering incidents & RQ3 \\
\bottomrule
\end{tabularx}
\end{table}

\subsection{Comparators}
\label{sec:design:comparators}

\paragraph{Thresholds (H1a).} Nine alternatives are evaluated on the same replay with their specified calibration and history settings: fixed thresholds from the first batch and the pooled warm-up; rolling mean and standard deviation; rolling median and median absolute deviation~\cite{Leys2013}; rolling Tukey fences~\cite{Tukey1977}; calibrated versions of both robust rules; exponentially weighted batch moments with $\lambda=0.3$~\cite{Roberts1959}; and the proposed formula using the previous rather than current batch. Pooled fixed thresholds are the primary comparator. Robust multipliers are calibrated to the warm-up flag rate of that comparator because the textbook multiplier of three flagged 27--29\% of valid records on skewed development columns. Two decomposition variants exchange numerical and category policies to identify which policy produces a difference.

Deequ~\cite{Schelter2018} is evaluated with default rule suggestions on observed warm-up and in a favourable configuration using extended suggestions on warm-up with injected errors removed. The latter excludes numerical suggestions that require exact equality with warm-up statistics. Deequ verifies its own constraints on each batch and supplies record violations. These configurations distinguish practical default behaviour from the benefit of additional calibration information.

\paragraph{Batch labels (H1b).} Comparators include the same gated drift signal with fixed, rolling robust or EWMA thresholds; the proposed classifier without the gate; a minimum-effect-size drift variant; static rules without drift; an AVH-inspired historical-statistics baseline~\cite{Tu2023}; Deequ anomaly detection as documented and warm-up calibrated; and TFDV~\cite{Caveness2020} using documented or calibrated thresholds and previous-batch or warm-up references. The AVH-inspired baseline uses Student prediction bounds and a nominal 5\% batch false-alarm budget, updating from all history or only alarm-free batches. It omits the published AVH constraint-selection optimiser. Its results therefore describe this restricted baseline, not a reproduction of the full AVH system.

\paragraph{Rules (H2a, H2b).} Fixed and adaptive statistical rules are evaluated alone and with proposed or admitted LLM rules. The dictionary-only model-authored comparison rules are tested before and after admission, and both candidate sources are evaluated alone. This separates the value of candidate generation, historical screening and combination with statistical rules. Category-guard policies follow Section~\ref{sec:method:guard}.

\paragraph{Schedules (H3).} Risk-based selection is compared with all-check execution, cost-matched periodic and random schedules, and failure-only or drift-only risk scores. Matching uses the risk run's realised standard-check execution rate. It is retrospective experimental matching, not an online budget-setting procedure; random execution counts can differ in a finite run, and matched check counts do not imply matched bytes or time.

\paragraph{Ablation.} The full factorial design crosses adaptation (pooled fixed or adaptive), LLM rules (absent or admitted) and selection (all or risk-based) on A--D, S1 and S2, giving eight configurations. The per-value category guard is evaluated separately and is not a factor in this design. Factorial results therefore do not establish the performance of an integrated guarded adaptive system.

\subsection{Measures}

For record detection, let TP, FP and FN denote true positives, false positives and false negatives, and TN denote true negatives. Precision is $\mathrm{TP}/(\mathrm{TP}+\mathrm{FP})$, recall is $\mathrm{TP}/(\mathrm{TP}+\mathrm{FN})$, and FPR is $\mathrm{FP}/(\mathrm{FP}+\mathrm{TN})$. F1 is the harmonic mean of precision and recall. Counts are pooled over evaluated batches and seeds for record-level micro-averages; undefined quantities, including recall in the zero-error control, are not interpreted as detection outcomes.

Recall on the injected domain-impossible subset is reported separately as \emph{semantic recall}. Batch evaluation uses four-class macro-F1, balanced accuracy, accuracy, pure-drift batches mislabelled degradation, and onset-to-detection delay in batches. The summary macro-F1 is calculated from pooled confusion counts; paired uncertainty uses per-seed macro-F1, a related but distinct estimand.

For candidate rules, measures include admission outcomes, the share of admitted rules exceeding 0.5\% FPR on valid evaluation records, and generation agreement on up to 200 sampled values per column plus synthetic probes. Category guards are assessed by valid and invalid value-admission rates and admission delay. Scheduling uses checks saved, estimated logical bytes touched and recall after onset. For sustained degradation (C/D), recall concerns all erroneous post-onset records; for H1--H4, it concerns incident-injected records during the incident. The H3 usefulness rule thresholds the reported mean recall rather than its confidence bound.

Timing uses five repetitions on the first test seed with randomised configuration order. Each repetition sums time over fifteen evaluated batches; the reported per-batch statistic is the median across repetitions of these totals divided by fifteen. Reuse applies to profile statistics and masks used in execution and drift gating. Logical bytes are implementation accounting estimates, not measured physical I/O or cloud charges.

\subsection{Statistical Analysis}

The replication unit is a simulation run, identified by its seed. Batches within a run share history, so individual records and batches are not treated as independent replications. Methods use the same seeds, and comparisons are paired. A cluster bootstrap resamples whole seeds with 5,000 replicates~\cite{FieldWelsh2007}. Counts are pooled within each resample for ratios such as recall and FPR, matching their point estimands. Seed reuse of a common source pool means uncertainty is conditional on these datasets and the specified replay process.

H1a is supported in a dataset--scenario case when the upper 95\% bound of adaptive-minus-fixed FPR is below zero and the lower recall-difference bound is above $-1$ pp. Only the FPR condition applies in Z. Decisions remain per dataset and scenario; they are not a study-wide familywise superiority claim.

The original RQ1a runner applied Holm adjustment~\cite{Holm1979} inside separate comparator calls, leaving one comparison per correction family. A separately documented reporting correction recomputes Holm across the sixteen comparator configurations within each dataset--scenario--metric family, without changing primary confidence intervals or H1a decisions. It is retrospective reporting repair, not a new frozen experiment. Bootstrap sign-tail $p$-values remain approximate, and zero among 5,000 resamples does not mean an exact probability of zero. The main interpretation therefore emphasises paired intervals, effect sizes and prespecified criteria. Non-significance is not evidence of equivalence.

\subsection{External Benchmarks}
\label{sec:design:external}

Tax and Beers from the Raha benchmark collection~\cite{Mahdavi2019} provide paired dirty and clean tables. Tax contains 200,000 synthetic tax records, including errors in combinations of individually plausible fields. Beers contains 2,410 US craft-beer records. Both are replayed as twenty shuffled batches per seed. Their warm-up contains uncontrolled source errors. Every evaluated Beers record is erroneous, whereas Tax has an erroneous-record share of about 1.46\%; thus the mostly-clean-record assumption fails much more strongly for Beers.

Ten percent of each table, with a minimum of 500 rows, is reserved for LLM prompts and excluded from replay. Ground-truth comparison normalises whitespace and numerically equivalent strings, so some formatting-only numerical differences are not labelled errors. Cell metrics are calculated over columns for which the monitor supplies checks; the share of all error cells represented by that denominator is stated with the results.

Tax also tests cross-column rules: dependencies such as zip code determining city, and conditional relations such as no child exemption when there are no children. Their structures are learned from warm-up or proposed by an LLM, with mappings learned from warm-up. Admission shares the 0.5\% warm-up flag-rate ceiling but uses mapping coverage and duplicate checks rather than the complete single-column algorithm. The external single-column runner uses empirical failure-mask redundancy, whereas the main coverage analysis reports logical redundancy screening. These configurations are distinguished when interpreting transfer. Beers uses cell metrics because record FPR is undefined and record recall is uninformative when every record is erroneous.

\subsection{Chronological Replay}
\label{sec:design:replay}

Raw Chicago inspections are replayed in natural order as monthly batches, with all checks executed. January--May 2010 provides warm-up; the incomplete final month is dropped. The evaluated period contains 196 months, June 2010--September 2026, and 309,502 records. Drift is compared with the previous five months and separately with the original warm-up, allowing long-term reference obsolescence to be distinguished from recent change.

The replay includes a documented procedure change on 1 July 2018. Around this transition, the share of \emph{Pass} results falls from 56\% to 12\%, while \emph{Pass w/ Conditions} rises from 12\% to 54\%. This is a known legitimate transition against which to inspect the operational label; one such event does not provide exhaustive drift ground truth.

Alarms are judged under explicit schema and domain policies where possible: missing values, exact duplicates, city-polygon membership and documented risk/result domains. Domain membership establishes admissibility, not factual accuracy of an individual inspection. Facility Type has no closed documented domain and requires independent human judgement (Section~\ref{sec:res:replay}). Check-record alarm counts can overlap and are not automatically counts of distinct erroneous rows.

\subsection{Implementation}

The monitor and harness use Python 3.14, pandas 3.0 and NumPy 2.5. Deequ runs as a Spark sidecar on OpenJDK 17 and TFDV in a container, on the same batches. The recorded experiments ran on one Apple Silicon laptop. The confirmatory run took 113 minutes with four parallel workers; sensitivity analysis took a further 79 minutes. This execution setting supports reproducibility of the reported implementation rather than a cloud-scale performance claim.

LLM responses are cached by fingerprints of prompts, data, settings and code. Thirty controlled-experiment generations (two models, three datasets, five generations) have an estimated total cost of US\$2.14 at configured prices. The cost calculation replaces implausible secondary-model input-token counts with a text-length estimate. This is an estimated generation expense, not audited billing, and excludes the separate per-value category queries.
